\section*{Introduction}

The masticatory system operates under dual control: a brainstem central pattern
generator (CPG) that produces the fundamental jaw rhythm~\cite{Lund1991,Grillner2006},
and a continuous stream of peripheral somatosensory signals from periodontal
mechanoreceptors, mucosal receptors, and muscle spindles that adjust the
amplitude and coordination of jaw-closing muscles on a cycle-by-cycle
basis~\cite{Trulsson2006,Lobbezoo1997,Sessle2006}.
Periodontal afferents encode bite force, tooth contact area, and food texture
with high spatial acuity and project via the trigeminal nuclear complex to the
masticatory motor nucleus, where they exert modulatory effects on jaw-closing
motor neuron pools~\cite{Turker2002,Sessle2006}.
The prevailing hierarchical model holds that the CPG is responsible for the
temporal patterning of jaw movement while peripheral feedback regulates the
gain of individual muscle contractions to maintain effective bite force against
varying mechanical loads~\cite{Lund1991,Trulsson2006}.
A core prediction of this model is that transient disruption of peripheral
afference should alter the amplitude of masticatory contractions while
leaving the fundamental rhythm intact.

A recent study from this laboratory~\cite{espinoza2025} examined the effects
of rhythmic chewing on working memory performance and frontocentral theta
oscillations using a 2\,$\times$\,2 within-subject crossover design (chewing
$\times$ anesthesia) in which unilateral topical lidocaine was compared against
a taste- and procedure-matched placebo.
That study demonstrated robust cognitive and electroencephalographic benefits
of chewing through a mechanism that appeared independent of oral somatosensory
feedback: neither behavioral nor EEG outcomes differed between anesthesia and
placebo conditions, leading to the interpretation that the cognitive effect
reflected a central mechanism.
The anesthesia factor was accordingly collapsed in the main analysis.
Critically, however, the surface masseter electromyogram (EMG) was not
analyzed as a primary outcome variable: prior analyses applied unpaired
nonparametric comparisons that are inappropriate for within-subject crossover
data~\cite{Cohen1988}, and the time-resolved structure of the pharmacological
manipulation -- topical lidocaine washout over approximately 15--20~minutes --
was not exploited.
Whether the lidocaine treatment produced any detectable physiological effect at
the motor level therefore remained an open question.

The present study addresses this gap with a secondary analysis of the same
dataset, focusing exclusively on the masseter EMG as a marker of
sensorimotor processing.
Topical lidocaine acts by reversibly blocking voltage-gated sodium channels,
reducing afferent discharge from the oral mucosa and superficial periodontal
tissues without crossing the blood-brain barrier and without directly
affecting the motor axons innervating the masseter~\cite{vonGunten2019}.
Spray application reliably anaesthetises mucosal surfaces; penetration to
deep periodontal ligament mechanoreceptors (located 0.5--3~mm below the
epithelial surface) is possible but less certain than with infiltration
injection, and is acknowledged as a limitation in the Discussion.
This pharmacological specificity nonetheless permits the interpretation that
any EMG change reflects altered oral sensory afference rather than direct
motor effects.
The crossover design with matched placebo controlled for non-pharmacological
components of the application procedure (taste, mucosal sensation, expectation).
The four chewing bouts per condition -- spaced at approximately 4.5-minute
intervals after a single spray application -- provide a natural time series
for tracking pharmacological washout.
We predicted that, consistent with the CPG gain-modulation framework,
lidocaine would increase masseter amplitude without altering chewing
frequency, and that this effect would decay as the anesthetic cleared.

The aims of this study were: (1) to determine whether unilateral
topical dental anesthesia produces a time-resolved change in masseter EMG;
(2) to characterise the dissociation between amplitude (gain) and timing
metrics (frequency, rate) as evidence for CPG-level robustness;
(3) to provide a post-hoc manipulation check for the parent study by
demonstrating that the anesthetic did produce a detectable physiological effect
at the motor periphery, thereby reinterpreting the absence of behavioral and
EEG anesthesia effects as a consequence of motor homeostasis rather than an
ineffective manipulation; and (4) to explore, as a secondary aim, whether
bilateral masseter synchrony was preserved under unilateral sensory reduction,
as an additional test of CPG robustness.

